\section{Buscas Não Informadas}
\label{sec:busca-nao-informada}

Nesta etapa foram implementados dois métodos de busca não informada:
a Busca em Profundidade (\textit{Depth-First Search} ou DFS) e a Busca
em Largura (\textit{Breadth-First Search} ou BFS). Esses algoritmos exploram o
espaço de estados utilizando apenas a estrutura de adjacência do
grafo, sem levar em consideração o custo das arestas ou estimativas
heurísticas sobre a proximidade do nó ao objetivo.

\subsection{Busca em Profundidade}
\label{subsec:dfs}

A Busca em Profundidade busca explorar o nó mais profundo (o mais longe do nó inicial)
presente na fronteira de busca. Caso atinja um estado sem sucessores válidos não visitados, o
algoritmo realiza o retorna para explorar ramos ainda não explorados.

Para implementar a política LIFO (\textit{Last-In, First-Out}) característica deste
algoritmo, a fronteira foi estruturada por meio de uma pilha (\texttt{list} em Python).
Cada entrada na pilha armazena uma tupla contendo o estado atual, a sequência de estados
que compõem o caminho percorrido desde a raiz e o custo acumulado da trajetória.

Como o espaço de estados do problema da Ponte e da Tocha contém múltiplos ciclos
(devido à possibilidade de pessoas realizarem travessias repetidas de ida e volta),
a DFS pura em árvore entraria em laços infinitos. Para garantir a terminação, foi
adotada a busca em grafo com um conjunto de estados visitados (\texttt{visitados}).
Um estado é formalmente adicionado a esse conjunto quando é retirado do topo da pilha
para expansão. Caso um estado retirado já pertença a \texttt{visitados}, ele é descartado.

A cada iteração, verifica-se a condição de parada através do método
\texttt{is\_estado\_final()}. Ao encontrar a configuração em que todas as pessoas e a
tocha estão no destino, a busca é finalizada, retornando o
caminho percorrido, o custo total associado e a quantidade de nós visitados. Além disso,
a rotina aceita um parâmetro opcional \texttt{limite\_nos}, que encerra a exploração
caso um número predeterminado de nós expandidos seja ultrapassado, protegendo a execução
contra buscas excessivamente longas.

Embora a DFS possua complexidade de espaço vantajosa em árvores gerais, em problemas
ponderados ela não oferece qualquer garantia de otimalidade. Como ela explora o
primeiro ramo disponível até o nó folha, a solução encontrada frequentemente inclui
movimentações redundantes e travessias demoradas, resultando em um tempo total de
travessia muito superior ao mínimo possível.

\subsection{Busca em Largura}
\label{subsec:bfs}

A Busca em Largura explora o grafo em camadas de profundidade, expandindo
todos os nós a uma distância de $d$ transições a partir do estado inicial antes de
avaliar qualquer nó a distância $d+1$.

A política FIFO (\textit{First-In, First-Out}) foi implementada utilizando uma fila
duplamente encadeada (\texttt{collections.deque}). Diferentemente da DFS, na BFS o
registro no conjunto de visitados ocorre no momento em que os nós sucessores são gerados
e enfileirados. Essa estratégia evita a inserção duplicada de estados idênticos em um
mesmo nível na fronteira, economizando memória de forma substancial.

A condição de parada é testada no momento em que o nó é desenfileirado. Se o estado
corresponder à meta, o algoritmo retorna a solução. Pela propriedade fundamental da BFS,
o primeiro caminho encontrado até o objetivo é garantidamente aquele com o
\textbf{menor número de ações (nós visitados)}, mas não necessariamente o com menor custo.
Como as travessias possuem custos diferentes (ditados pela pessoa
mais lenta do movimento), um caminho com menor número de passos pode apresentar um custo
acumulado significativamente mais alto do que um caminho alternativo com o mesmo número
de passos ou passos adicionais que sincronize melhor os tempos dos indivíduos.
